% =====================================================================
%  Muc VII viet lai -- TETC-2026-05-0252
%
%  Ban da nop hua ba thu la "ongoing work" hoac "deferred": (a) do tac dong
%  cua be rong mach, (b) so voi baseline khong phai SVM, (c) dataset thu hai.
%  Ca ba nay bay gio deu co so do, nen muc han che khong con la loi hua ma la
%  bao cao ket qua kem ranh gioi ro rang.
%
%  Cho DUY NHAT van con la loi hua: thuc thi tren QPU that. Phai ghi that ro
%  da lam gi va CHUA lam gi -- neu de mo ho thi reviewer se doc "NISQ-aware"
%  thanh mot khang dinh ve phan cung ma bai khong do duoc.
% =====================================================================

\section{Limitations and Threats to Validity}
\label{sec:limitations}

\subsection{Hardware noise: what was and was not measured}
\revnote{rewrite}{The submitted version covered finite-shot error only. A backend-derived noise model is added (AE-6, R1-6, R3-4), with a statement of what is not measured.}
\begin{revbody}

The submitted version reported only finite-shot sampling error and deferred
coherent noise to a companion study. Reviewers correctly objected that shot
noise alone cannot support a NISQ-feasibility claim. We therefore state the
scope precisely.

\emph{What we measure.} The C2 pipeline is re-run under a noise model derived
from the published calibration data of a real superconducting backend
(\texttt{FakeManilaV2}) \cite{ref27} via \texttt{NoiseModel.from\_backend}, which
carries gate errors, readout errors, and gate-duration-dependent thermal
relaxation.
Circuits are transpiled onto that backend's coupling map and basis gate set,
giving a depth of $59$ and $44$ two-qubit gates for the $n=4$, $r=2$ map.
Under this model the entanglement conclusion is unchanged: the alignment gap
between QSVM-ZZ and QSVM-Z survives, while the macro-$F_1$ difference remains
within its confidence interval. This noisy re-run covers one of the ten runs
and the $300$-record test split (Table~\ref{tab:protocol}); it establishes
that the noise model does not move the operating point, not a noisy confidence
interval, which would require the remaining nine runs at the same $512$ shots
per kernel entry.

\emph{What we do not measure.} The circuits were not executed on a quantum
processor. Two effects are therefore outside our evidence: drift of calibration
parameters between the snapshot and execution time, and idle-time decoherence
beyond what gate durations already attribute. Our released code contains the
execution path for a real backend, and the comparison it would produce
(exact statevector $\rightarrow$ backend-derived noise $\rightarrow$ hardware)
is specified, but we report no hardware numbers and make no claim about them.
Readers should read ``NISQ-aware'' as a statement about \emph{circuit budget}
(qubit count, depth, and two-qubit gate count chosen under an explicit cost
model, in the sense of \cite{ref8}) and not as a claim of hardware validation.
\end{revbody}

\subsection{Embedding size: now measured, and the news is not good for width}
\revnote{fix}{The submitted version left this as an ``open question''; it is now measured, and the result is unfavourable to wider circuits.}
\begin{revbody}

The submitted version listed the effect of moving along the frontier as an open
question. We have since measured it, and the answer constrains the method more
sharply than we expected.

Increasing the feature budget $K$ does buy accuracy for a classical proxy, but
it is not free: retaining $85\%$ of the variance at $K=80$ requires $n^\ast=8$
qubits rather than the $n^\ast=4$ that $K=20$ admits, raising the two-qubit gate
count from $24$ to $112$. Re-running the full sample-complexity sweep at that
operating point removes the quantum kernel's advantage entirely: across eight
training-set sizes and six baselines, all $48$ paired comparisons favour the
classical model, and the entanglement ablation reverses sign, with QSVM-ZZ now
\emph{below} its entanglement-free control.

The mechanism is measurable rather than conjectural. The dispersion of the
off-diagonal Gram entries decays with width as $n^{-\alpha}$, and the fitted
exponent for the ZZ map is close to twice that of the Z map at $K=20$
($0.59$ vs.\ $0.29$), consistent with the $\binom{n}{2}$ pairwise phase terms of
the former growing quadratically where the latter's grow linearly. Across
$n=4$ to $10$ the macro-$F_1$ of the ZZ kernel tracks this dispersion
($r=+0.77$, 95\% CI $[+0.04,+0.96]$), whereas the Z kernel, which barely
concentrates, shows no resolved dependence ($r=+0.32$, $[-0.57,+0.87]$). This is
the concentration phenomenon of \cite{thanasilp2024} appearing directly in our
setting.

Three caveats. The exponent ratio is $2.02$ at $K=20$, where the kernel does not
collapse, and $3.32$ at $K=80$, where it does, so the term-counting argument
predicts the ordering of the two rates and, at the collapse, not their ratio;
the abstract and Sec.~\ref{sec:intro} state both values. The $F_1$--dispersion
relation is measured on seven widths at a single training-set size, and the
contrast between the two kernels' correlations is not resolvable at seven
points ($z=0.97$, $p=0.33$); we present the relation as an empirical regularity
of the \ZZ{} map, not a law, and make no claim that it is specific to the
concentrating map. A test of either would need more widths and more than one
training-set size.
\end{revbody}

\subsection{Classical-baseline strength}
\revnote{rewrite}{RF and XGBoost added; CatBoost, TabNet and FT-Transformer deliberately not added, with the reason stated (R1-5).}
\begin{revbody}

The submitted version compared only against SVM variants. We have added a
random forest and XGBoost under the identical protocol, and they change the
headline: on the main benchmark XGBoost attains the highest mean macro-$F_1$
($0.8503$) ahead of QSVM-ZZ ($0.8469$), with overlapping intervals. We report
this as an ordering of point estimates, not a significant separation.

We did not add CatBoost or deep tabular models such as TabNet and
FT-Transformer. CatBoost would supply a second gradient-boosting representative
without changing the comparison structure; the deep tabular family would require
its own tuning-budget protocol to be compared fairly, which would enlarge the
study from a controlled quantum-kernel benchmark into a survey of tabular
architectures. We therefore do not claim coverage of all strong tabular methods,
and any practical reading of our results should be qualified accordingly.
\end{revbody}

\subsection{Dataset breadth and protocol}
\revnote{rewrite}{UNSW-NB15 added, with duplication statistics for both corpora.}
\begin{revbody}

The submitted version rested almost entirely on NSL-KDD. We have added
UNSW-NB15 under the same protocol, which is what allows us to separate findings
that transfer from findings that do not: the dimension-selection rule and the
entanglement ablation transfer, while the sample-complexity crossover and the
rare-class behaviour do not.

Two properties of the data limit what either dataset can show. NSL-KDD is an
aging benchmark in which $610$ test rows are identical to a training row in both
features and label ($617$ match on features alone). UNSW-NB15 is far more heavily
duplicated: $47.3\%$ of training rows are internal repeats, and $25.0\%$ of test
rows have an exact duplicate in the training partition, which
inflates absolute scores. Our claims are paired differences between models
evaluated on the same rows, which removes the duplication from the level but
not necessarily from the comparison: a tree ensemble reproduces an exact
training row by construction, a fidelity-kernel SVM does so only through its
support vectors and regulariser, and on the quarter of the UNSW-NB15 test set
that duplicates a training row a paired difference may therefore measure
memorisation rather than generalisation. We have not re-scored the paired
differences on the de-duplicated test subsets, because the sweep's per-row
predictions were not retained, and we list that re-scoring as the one cheap
check this study owes: until it is done, the UNSW-NB15 verdicts, and to a
smaller degree the NSL-KDD ones, carry this threat, and the absolute numbers
are not comparable to studies that de-duplicate.

Finally, the feature budget $K=20$ is inherited from the submitted protocol and
held fixed so that the revised results remain comparable to it.
Fig.~\ref{fig:ksweep} shows it is not the accuracy optimum; it is the value at
which the selection rule returns a four-qubit circuit.
\end{revbody}
